\section*{Capítulo \ref{ch:g6:classification-kinship} --- Classificação e parentesco}

\begin{solution}{exo:g6:classification-kinship:1}
Uma característica que um organismo possui, enunciada de forma
conferível: tem \omterm{def:g4:classifying-animals:vertebrate}{coluna vertebral}; tem quatro \omterm{def:g1:my-body:parts}{membros}; tem pelo e
alimenta os filhotes com \omterm{def:g2:animals-grow-up:born}{leite} (ou: tem \omterm{def:g1:animals-around-us:covering}{penas}).
\end{solution}

\begin{solution}{exo:g6:classification-kinship:2}
O conjunto de donos de cada \omterm{def:g6:classification-kinship:attribute}{atributo} fica inteiramente dentro do
conjunto de outro: todos os que têm pelo têm quatro \omterm{def:g1:my-body:parts}{membros}, todos os
\omterm{def:g1:animals-around-us:animal}{animais} de quatro \omterm{def:g1:my-body:parts}{membros} têm \omterm{def:g4:classifying-animals:vertebrate}{coluna vertebral} --- nunca o contrário,
nunca cruzando. Caixas dentro de caixas, sem ninguém montado no muro.
\end{solution}

\begin{solution}{exo:g6:classification-kinship:3}
Porque elas os herdaram de um ancestral comum: um \omterm{def:g6:classification-kinship:attribute}{atributo} desce por
uma linhagem de família, então os donos dele são os descendentes
daquele ancestral --- uma família.
\end{solution}

\begin{solution}{exo:g6:classification-kinship:4}
A truta está além só do pontinho da \omterm{def:g4:classifying-animals:vertebrate}{coluna vertebral}. A baleia está
além da \omterm{def:g4:classifying-animals:vertebrate}{coluna vertebral}, dos quatro \omterm{def:g1:my-body:parts}{membros} e do pelo e \omterm{def:g2:animals-grow-up:born}{leite} --- dos
três.
\end{solution}

\begin{solution}{exo:g6:classification-kinship:5}
A baleia e o gato: eles compartilham a bifurcação do pelo e \omterm{def:g2:animals-grow-up:born}{leite}, a
mais nova desenhada; a truta se separou na bifurcação mais antiga.
\end{solution}

\begin{solution}{exo:g6:classification-kinship:6}
São a mesma informação duas vezes: uma caixa é um ramo com todos os
galhinhos dele --- tudo o que está além de uma bifurcação herda o
\omterm{def:g6:classification-kinship:attribute}{atributo} dela e fica na caixa dela.
\end{solution}

\begin{solution}{exo:g6:classification-kinship:7}
Pelo e \omterm{def:g2:animals-grow-up:born}{leite} são herdados dentro de uma única linhagem de família ---
isso marca ancestralidade. ``Tem formato de peixe'' pode ser carimbado
em estranhos pelo mesmo modo de vida aquático, como a baleia mostra:
formato emprestado, família em outro lugar. A \omterm{def:g4:classifying-animals:classification}{classificação} separa
famílias, então confia na herança e desconfia do casaco emprestado.
\end{solution}

\begin{solution}{exo:g6:classification-kinship:8}
Truta: só \omterm{def:g4:classifying-animals:vertebrate}{coluna vertebral}. Sapo: \omterm{def:g4:classifying-animals:vertebrate}{coluna vertebral}, quatro \omterm{def:g1:my-body:parts}{membros}.
Gato: \omterm{def:g4:classifying-animals:vertebrate}{coluna vertebral}, quatro \omterm{def:g1:my-body:parts}{membros}, pelo e \omterm{def:g2:animals-grow-up:born}{leite}. Sabiá: \omterm{def:g4:classifying-animals:vertebrate}{coluna
vertebral}, quatro \omterm{def:g1:my-body:parts}{membros}, \omterm{def:g1:animals-around-us:covering}{penas}. O encaixe aparece nos vistos: os
vistos de toda linha incluem a \omterm{def:g4:classifying-animals:vertebrate}{coluna vertebral}; toda linha com visto
de revestimento também tem visto de quatro \omterm{def:g1:my-body:parts}{membros}.
\end{solution}

\begin{solution}{exo:g6:classification-kinship:9}
Ele passa pela \omterm{def:g4:classifying-animals:vertebrate}{coluna vertebral} e pelos quatro \omterm{def:g1:my-body:parts}{membros} e sai num ramo
próprio, com a \omterm{def:g1:the-five-senses:sense}{pele} seca e escamosa dele --- ao lado, e não além, das
bifurcações do pelo e das \omterm{def:g1:animals-around-us:covering}{penas}.
\end{solution}

\begin{solution}{exo:g6:classification-kinship:10}
Amamentar marca \omterm{prop:g6:classification-kinship:kinship}{parentesco}: pelo e \omterm{def:g2:animals-grow-up:born}{leite} são a herança da família dos
\omterm{prop:g3:sorting-living-things:groups}{mamíferos}, então o morcego se ramifica junto com o gato. Voar é
emprestado pelo modo de vida --- as asas surgiram separadamente no ramo
das \omterm{prop:g3:sorting-living-things:groups}{aves} e no ramo dos morcegos, como o corpo hidrodinâmico surgiu na
baleia e na truta.
\end{solution}

\begin{solution}{exo:g6:classification-kinship:11}
Ele pertence. O ornitorrinco descende do ancestral comum dos \omterm{prop:g3:sorting-living-things:groups}{mamíferos}
e carrega as heranças da família (pelo, \omterm{def:g2:animals-grow-up:born}{leite}); o \omterm{def:g2:animals-grow-up:born}{ovo} dele é uma
característica antiga que o ramo primitivo dele manteve enquanto os
outros a perderam. A família é o ancestral e \emph{todos} os
descendentes --- esquisitices incluídas; o pacote
(\cref{exo:g4:classifying-animals:11}) já o tinha aprovado.
\end{solution}

\begin{solution}{exo:g6:classification-kinship:12}
Se todo par de ramos se encontra, então as \omterm{def:g5:biodiversity:species}{espécies} precisam mudar ao
longo das gerações e de muito tempo --- com \omterm{def:g6:classification-kinship:attribute}{atributos} novos surgindo e
sendo herdados --- para que uma linhagem ancestral pudesse se ramificar
em truta, gato e margarida. A implicação numa frase: todos os \omterm{def:g6:exploring-our-environment:components}{seres
vivos} descendem, com mudança, de ancestrais comuns. Provas a exigir:
vestígios da mudança --- formas intermediárias, restos preservados em
rochas antigas e semelhanças de família situadas exatamente onde a
\omterm{def:g6:classification-kinship:tree}{árvore} prevê (assunto de um ano posterior).
\end{solution}

\begin{solution}{pb:g6:classification-kinship:1}
\textbf{1.} Truta: \omterm{def:g4:classifying-animals:vertebrate}{coluna vertebral}. Sapo: \omterm{def:g4:classifying-animals:vertebrate}{coluna vertebral}, quatro
\omterm{def:g1:my-body:parts}{membros}. Lagarto: \omterm{def:g4:classifying-animals:vertebrate}{coluna vertebral}, quatro \omterm{def:g1:my-body:parts}{membros}, \omterm{def:g1:the-five-senses:sense}{pele} seca e
escamosa. Pombo: \omterm{def:g4:classifying-animals:vertebrate}{coluna vertebral}, quatro \omterm{def:g1:my-body:parts}{membros}, \omterm{def:g1:animals-around-us:covering}{penas}. Camundongo:
\omterm{def:g4:classifying-animals:vertebrate}{coluna vertebral}, quatro \omterm{def:g1:my-body:parts}{membros}, pelo e \omterm{def:g2:animals-grow-up:born}{leite}. Golfinho: \omterm{def:g4:classifying-animals:vertebrate}{coluna
vertebral}, quatro \omterm{def:g1:my-body:parts}{membros}, pelo e \omterm{def:g2:animals-grow-up:born}{leite} (com o testemunho da etiqueta
sobre o \omterm{def:g2:animals-grow-up:born}{leite}; quem decide são as características, e o \omterm{def:g2:animals-grow-up:born}{leite} é a
característica decisiva).

\textbf{2.} A \omterm{def:g4:classifying-animals:vertebrate}{coluna vertebral} --- os seis são \omterm{def:g4:classifying-animals:vertebrate}{vertebrados}. O
compartilhamento mais amplo marca o ancestral comum mais antigo: os
seis são uma família só naquela profundidade.

\textbf{3.} Sapo, lagarto, pombo, camundongo, golfinho. As asas do
pombo e as nadadeiras do golfinho são quatro \omterm{def:g1:my-body:parts}{membros} em roupa de
trabalho: os mesmos \omterm{def:g1:my-body:parts}{membros} herdados, moldados por vidas diferentes ---
a \omterm{def:g6:classification-kinship:tree}{árvore} conta a herança, e não a roupa.

\textbf{4.} Porque um endereço compartilhado não é uma família
compartilhada: ``vive na água'' é um \omterm{def:g2:where-animals-live:habitat}{hábitat}, que pode ser emprestado
por estranhos, e não um \omterm{def:g6:classification-kinship:attribute}{atributo} da ancestralidade do corpo --- a velha
regra da \cref{def:g3:sorting-living-things:criterion}, agora com o
motivo dela.

\textbf{5.} Tronco $\rightarrow$ bifurcação da \omterm{def:g4:classifying-animals:vertebrate}{coluna vertebral} (a
truta sai depois dela) $\rightarrow$ bifurcação dos quatro \omterm{def:g1:my-body:parts}{membros}
$\rightarrow$ três ramos de revestimento: escamas secas (lagarto),
\omterm{def:g1:animals-around-us:covering}{penas} (pombo), pelo e \omterm{def:g2:animals-grow-up:born}{leite} (camundongo e golfinho juntos na
bifurcação mais nova); o sapo sai depois dos quatro \omterm{def:g1:my-body:parts}{membros}, antes de
todo revestimento.

\textbf{6.} O primo mais próximo do camundongo na gaveta é o golfinho
(eles compartilham a bifurcação do pelo e \omterm{def:g2:animals-grow-up:born}{leite}). O da truta são todos
igualmente --- ela se separou na primeira bifurcação, então os outros
cinco são primos igualmente distantes dela.

\textbf{7.} Um anfíbio --- de quatro \omterm{def:g1:my-body:parts}{membros}, \omterm{def:g1:the-five-senses:sense}{pele} nua e úmida, sem
bifurcação de revestimento própria; o ramo dele sai entre a bifurcação
dos quatro \omterm{def:g1:my-body:parts}{membros} e os revestimentos.

\textbf{8.} Além da \omterm{def:g4:classifying-animals:vertebrate}{coluna vertebral} e dos quatro \omterm{def:g1:my-body:parts}{membros}, no ramo das
escamas secas: um réptil --- classificado sem ser visto, porque quem
decide são os \omterm{def:g6:classification-kinship:attribute}{atributos}, e não as etiquetas.

\textbf{9.} O golfinho e o camundongo ficam mais perto: pelo e \omterm{def:g2:animals-grow-up:born}{leite}
são uma bifurcação herdada. O corpo hidrodinâmico é o uniforme da
água, distribuído igualmente a nadadores sem \omterm{prop:g6:classification-kinship:kinship}{parentesco}
(\cref{exo:g4:adapted-to-their-world:11}); se fosse contado, ele
arquivaria roupa emprestada como família e arruinaria o encaixe.

\textbf{10.} Uma única \omterm{def:g5:biodiversity:species}{espécie} ancestral comum --- um \omterm{def:g4:classifying-animals:vertebrate}{vertebrado}
primitivo de quatro \omterm{def:g1:my-body:parts}{membros} --- da qual o sapo, o lagarto, o pombo, o
camundongo e o golfinho todos descendem, cada um carregando os quatro
\omterm{def:g1:my-body:parts}{membros} dele numa roupa diferente.

\textbf{11.} Que os seis --- e a turma que os desenhou --- pertencem a
uma única família que se ramifica, descendente com mudança de
ancestrais comuns, com os seres humanos no ramo do pelo e \omterm{def:g2:animals-grow-up:born}{leite}, entre
os outros.

\textbf{12.} O encaixe é um fato sobre o mundo vivo --- os \omterm{def:g6:classification-kinship:attribute}{atributos}
realmente caem em caixas dentro de caixas, e isso não foi decretado
por arquivista nenhum. A \omterm{def:g6:classification-kinship:tree}{árvore} explica o fato: a herança a partir de
ancestrais comuns \emph{precisa} produzir \omterm{def:g6:classification-kinship:attribute}{atributos} encaixados,
enquanto o mero arquivamento deixaria a arrumação obediente do mundo
como um milagre.
\end{solution}
